\documentclass[11pt,a4paper]{article}
\usepackage[spanish,es-noquoting]{babel}
\usepackage[T1]{fontenc}
\usepackage[utf8]{inputenc}
\usepackage{lmodern}
\usepackage{microtype}
\usepackage[a4paper,margin=2.35cm,headheight=14pt]{geometry}
\usepackage{amsmath,amssymb}
\usepackage{siunitx}
\sisetup{output-decimal-marker={,},per-mode=symbol}
\usepackage{circuitikz}
\usepackage{tikz}
\usepackage{booktabs,array,tabularx,longtable}
\usepackage{enumitem}
\usepackage{fancyhdr}
\usepackage{float}
\usepackage{needspace}
\usepackage{listings}
\usepackage[hidelinks]{hyperref}
\setlength{\parindent}{0pt}
\setlength{\parskip}{0.48em}
\setlist[itemize]{leftmargin=1.35em,itemsep=0.1em,topsep=0.2em}
\setlist[enumerate]{leftmargin=1.5em,itemsep=0.1em,topsep=0.2em}
\renewcommand{\arraystretch}{1.14}
\pagestyle{fancy}
\fancyhf{}
\fancyfoot[C]{\thepage}
\renewcommand{\headrulewidth}{0pt}
\renewcommand{\footrulewidth}{0pt}
\hypersetup{colorlinks=false,pdfborder={0 0 0}}
\lstset{basicstyle=\ttfamily\small,keywordstyle=\color{black},commentstyle=\color{black},stringstyle=\color{black},frame=single,framerule=0.5pt,columns=fullflexible,keepspaces=true,showstringspaces=false,breaklines=true,xleftmargin=0.3em,xrightmargin=0.3em,aboveskip=0.7em,belowskip=0.7em}
\newcommand{\titular}[2]{%
  {\LARGE\bfseries #1\par}
  \vspace{0.45em}
  {\normalsize #2\par}
  \vspace{0.45em}\hrule\vspace{0.9em}
}

\begin{document}
\titular{Detector de paso por infrarrojos}{Circuitos, cálculos y medidas del prototipo \hfill 06/11/2026}

\section{Objetivo y estado del trabajo}
Quiero construir una barrera óptica para detectar cuándo alguien cruza una puerta. El emisor utiliza un LED infrarrojo y el receptor es un fototransistor. El microcontrolador se encarga de activar el emisor y leer la tensión generada en el receptor.

En esta fase estoy utilizando mi placa ESP32-S3 N16R8 y una protoboard. Esta placa me sirve para comprobar la electrónica y preparar el programa. La PCB definitiva será un diseño propio con el microcontrolador soldado directamente, pero todavía no he elegido ese componente ni he implementado las notificaciones.

\begin{figure}[H]
\centering
\begin{tikzpicture}[>=latex, every node/.style={font=\small}]
\node[draw,minimum width=3.15cm,minimum height=1.02cm,align=center] (tx) at (0,0) {LED infrarrojo\\$L$-53F3C + BC337};
\node[draw,minimum width=3.1cm,minimum height=1.02cm,align=center] (rx) at (7,0) {Fototransistor\\BPW40};
\draw[->,dashed] (tx.east) -- node[above]{Luz infrarroja} (rx.west);
\node[draw,minimum width=3.2cm,minimum height=.86cm,align=center] (mcu) at (3.5,-2.15) {ESP32-S3\\GPIO4 / ADC GPIO5};
\draw[->] (mcu.west) -| (tx.south);
\draw[->] (rx.south) |- (mcu.east);
\end{tikzpicture}
\caption{Bloques utilizados para comprobar el detector. No es todavía un esquema de fabricación.}
\end{figure}

\subsection{Componentes de la prueba}
\begin{center}
\begin{tabularx}{\linewidth}{@{}l l X@{}}
\toprule
Referencia & Componente & Función\\
\midrule
U1 & ESP32-S3 N16R8 & Programación y lectura ADC\\
D1 & L-53F3C (940 nm) & Emisor infrarrojo\\
Q1 & BC337 & Conmutación del emisor\\
Q2 & BPW40 & Receptor óptico\\
R1 & $\SI{180}{\ohm}$ & Limitación de corriente del LED IR\\
R2 & $\SI{1}{\kilo\ohm}$ & Limitación de corriente de base\\
R3 & $\SI{10}{\kilo\ohm}$ & Pull-up del receptor\\
\bottomrule
\end{tabularx}
\end{center}

La alimentación de las pruebas procede de \texttt{3V3} de la placa, alimentada a su vez por USB-C. Todos los puntos de masa son comunes.

\section{Comprobación inicial de GPIO4}
Antes del emisor infrarrojo conecté un LED visible y una resistencia de $\SI{330}{\ohm}$ para comprobar la salida digital. El programa alternaba \texttt{HIGH} y \texttt{LOW}; después utilicé estados fijos para facilitar las medidas con el multímetro.

\begin{figure}[H]
\centering
\begin{circuitikz}[american currents, scale=0.95, transform shape]
\draw (0,0) node[left]{GPIO4} to[R,l=$\SI{330}{\ohm}$] (3,0) to[led,l=LED visible] (5,0) -- (5,-1) node[ground]{};
\end{circuitikz}
\caption{Primera prueba de salida digital.}
\end{figure}

Si tomo una caída típica de $\SI{2.0}{\volt}$ para un LED rojo, la corriente estimada es
\begin{equation}
I_{\mathrm{LED}} \simeq \frac{3{,}3-2{,}0}{330}
  =\SI{3.94}{\milli\ampere}.
\end{equation}
Este resultado es un cálculo previo, no una corriente medida. La prueba confirmó que GPIO4 podía conmutar aproximadamente entre $\SI{0}{\volt}$ y $\SI{3.3}{\volt}$.

\section{Emisor infrarrojo con BC337}
Para no hacer pasar toda la corriente del LED por el GPIO, añadí el BC337 como interruptor por el lado de masa. El LED se alimenta a través de R1 y el transistor recibe su señal de control mediante R2.

\begin{figure}[H]
\centering
\begin{circuitikz}[american currents, scale=1, transform shape]
\node[npn] (Q) at (4,1.0) {};
\draw (4,5.1) node[above]{3V3}
      to[R,l={$R_1: \SI{180}{\ohm}$}] (4,3.6)
      to[led,l={$D_1$ (IR)}] (4,2.3)
      -- (Q.C);
\draw (Q.E) -- (4,-.5) node[ground]{};
\draw (Q.B) to[R,l_={$R_2: \SI{1}{\kilo\ohm}$}] (0,1.0) node[left]{GPIO4};
\node[right] at (Q.C) {\hspace{.25cm}C};
\node[right] at (Q.E) {\hspace{.25cm}E};
\node[above left] at (Q.B) {B};
\node[right] at (4.75,1.0) {BC337};
\end{circuitikz}
\caption{Driver del LED infrarrojo. El ánodo del LED queda hacia R1 y el cátodo hacia el colector.}
\end{figure}

\subsection{Dimensionado}
Para estimar la corriente del emisor utilicé $V_F\simeq\SI{1.2}{\volt}$ y $V_{CE(\mathrm{sat})}\simeq\SI{0.2}{\volt}$. Son valores de cálculo, no medidas de esta unidad.
\begin{equation}
I_{D1}=\frac{V_{CC}-V_F-V_{CE(\mathrm{sat})}}{R_1}
=\frac{3{,}3-1{,}2-0{,}2}{180}
\simeq\SI{10.6}{\milli\ampere}.
\end{equation}

Para la base, tomando $V_{BE}\simeq\SI{0.7}{\volt}$:
\begin{equation}
I_B\simeq\frac{3{,}3-0{,}7}{1000}
=\SI{2.6}{\milli\ampere}.
\end{equation}

La disipación prevista de R1 con la corriente estimada queda aproximadamente en
\begin{equation}
P_{R1}=I_{D1}^{2}R_1\simeq\SI{0.020}{\watt}.
\end{equation}

\subsection{Medidas realizadas}
\begin{center}
\begin{tabularx}{\linewidth}{@{}X r X@{}}
\toprule
Magnitud & Valor & Observación\\
\midrule
$V_{GPIO4}$ con HIGH & $\simeq\SI{3.3}{\volt}$ & Estado alto comprobado\\
$V_{C}$ respecto a GND & $\SI{0.009}{\volt}$ & Medido con HIGH fijo\\
Caída en R1 & $\simeq\SI{2.16}{\volt}$ & Medida con el LED activado\\
$I_{D1}$ deducida de R1 & $\SI{12.0}{\milli\ampere}$ & Calculada, no medida directamente\\
\bottomrule
\end{tabularx}
\end{center}

A partir de la caída de tensión en R1:
\begin{equation}
I_{D1}=\frac{V_{R1}}{R_1}=\frac{2{,}16}{180}
=\SI{12.0}{\milli\ampere}.
\end{equation}

La tensión de colector de $\SI{0.009}{\volt}$ indica que el nodo estaba muy próximo a masa durante esa prueba. Como no tomé todas las medidas a la vez, no utilizaré esas lecturas para estimar con precisión la caída del LED. El emisor queda validado para continuar con las pruebas ópticas.

\textbf{Procedimiento de medida:} para medir R1 puse el multímetro en tensión continua, con una punta en cada extremo de la resistencia. Para medir directamente corriente en modo mA habría que abrir la rama e introducir el instrumento en serie.

\section{Receptor con fototransistor BPW40}
El receptor utiliza una resistencia de pull-up. La señal que leo en GPIO5 es la tensión del nodo entre R3 y el colector del fototransistor.

\begin{figure}[H]
\centering
\begin{circuitikz}[scale=1,transform shape]
\node[npn] (P) at (3.5,1.1) {};
\draw (3.5,4.8) node[above]{3V3}
  to[R,l={$R_3: \SI{10}{\kilo\ohm}$}] (3.5,3.0)
  -- (3.5,2.35) node[circ] (N) {}
  -- (P.C);
\draw (N) -- (6.2,2.35) node[right]{GPIO5 (ADC)};
\draw (P.E) -- (3.5,-.4) node[ground]{};
\node[right] at (4.2,1.1) {BPW40};
\draw[->] (1.55,2.6) -- (2.45,1.85);
\draw[->] (1.65,1.9) -- (2.5,1.2);
\node[left,align=right] at (1.6,2.4) {Luz IR};
\end{circuitikz}
\caption{Conexión eléctrica del receptor. La base dibujada en el símbolo de transistor no lleva una conexión externa.}
\end{figure}

Cuando el BPW40 recibe más luz, circula más corriente hacia masa y la tensión del nodo GPIO5 disminuye. Si lo tapo, el fototransistor conduce menos y R3 eleva la tensión hacia $\SI{3.3}{\volt}$.

Para una tensión medida $V_{GPIO5}$ puedo calcular la corriente por la resistencia pull-up con
\begin{equation}
I_{R3}=\frac{3{,}3-V_{GPIO5}}{\SI{10}{\kilo\ohm}}.
\end{equation}

El ESP32-S3 dispone de un ADC de 12 bits, con lecturas digitales de 0 a 4095. Esos valores no deben tratarse como voltios sin considerar la configuración y calibración del ADC.

\subsection{Primeras lecturas}
\begin{center}
\begin{tabularx}{\linewidth}{@{}X c X@{}}
\toprule
Condición & ADC & Interpretación\\
\midrule
Iluminación normal & 3940--3945 & Algo de luz llega al BPW40\\
Mano en la zona del haz & 4057--4067 & Disminuye la luz recibida\\
Receptor totalmente tapado & 4095 & Lectura en el máximo del ADC\\
\bottomrule
\end{tabularx}
\end{center}

La diferencia entre las dos primeras condiciones ronda 120 cuentas. No he comprobado todavía que esas variaciones se deban exclusivamente a mi LED IR: el sensor también responde a la iluminación de la habitación.

Al principio el ADC se quedaba fijo en 4095. El problema era que había conectado el BPW40 con la polaridad invertida. Al corregir colector y emisor, el valor empezó a responder a la luz.

\section{Montaje conjunto y posición de los sensores}

\begin{figure}[H]
\centering
\begin{tikzpicture}[>=latex, every node/.style={font=\small}]
\draw (0,0) -- (12.7,0);
\node[draw,minimum width=2.25cm,minimum height=0.9cm] (led) at (2,1.25) {LED IR};
\node[draw,minimum width=2.25cm,minimum height=0.9cm] (rx) at (10.7,1.25) {BPW40};
\draw[->,dashed,thick] (led.east) -- node[above]{Eje óptico} (rx.west);
\draw (1.7,.8) -- (1.7,0);
\draw (10.4,.8) -- (10.4,0);
\node[below] at (6.35,0) {Protoboard (representación lateral)};
\end{tikzpicture}
\caption{Para la prueba, las puntas de los dos componentes deben mirarse. Estar en la misma columna de la protoboard no garantiza una buena alineación óptica.}
\end{figure}

En la primera fotografía, el emisor y el receptor estaban apuntando hacia arriba. La corrección prevista consiste en orientar las dos cápsulas una frente a otra, empezando a 2--5 cm de distancia. Doblaré las patillas con cuidado, sin forzar la base del encapsulado ni permitir cortocircuitos entre terminales.

\section{Prueba diferencial y próximas medidas}
En lugar de mirar únicamente un valor absoluto del ADC, el código preparado mide el receptor con el emisor apagado y encendido. Calculo:
\begin{equation}
\Delta ADC=ADC_{\text{sin IR}}-ADC_{\text{con IR}}.
\end{equation}
Si nuestro emisor ilumina correctamente al receptor, espero un $\Delta ADC$ positivo. Al interrumpir el haz con un objeto opaco, la diferencia debería hacerse menor. Esta técnica ayuda a separar la contribución del LED de una iluminación ambiente que permanezca aproximadamente constante durante ambas lecturas.

\Needspace{14\baselineskip}
\begin{lstlisting}[language=C++,caption={Fragmento del ensayo diferencial preparado en PlatformIO.}]
digitalWrite(irEmisor, LOW);
delay(200);
int sinIR = medirReceptor();

digitalWrite(irEmisor, HIGH);
delay(200);
int conIR = medirReceptor();

int diferencia = sinIR - conIR;
\end{lstlisting}

\textbf{Esta prueba todavía está pendiente de validar} con los sensores orientados frente a frente. Tampoco se ha implementado una portadora de 38 kHz. Encender y apagar el LED lentamente no equivale a una barrera IR modulada.

\subsection{Tareas pendientes}
\begin{enumerate}
\item Comparar los valores \texttt{sinIR}, \texttt{conIR} y \texttt{diferencia} con el haz libre y bloqueado.
\item Repetir la medida con diferentes condiciones de luz ambiental y a mayor distancia.
\item Elegir el método de detección y comprobar falsos positivos y tiempo de respuesta.
\item Implementar un aviso al móvil y registrar las pruebas.
\item Seleccionar el microcontrolador desnudo, diseñar el esquemático en KiCad y fabricar la PCB.
\end{enumerate}

\section{Documentación utilizada}
\begin{itemize}
\item Kingbright, \emph{L-53F3C} (datasheet del LED emisor IR).
\item onsemi, \emph{BC337} (datasheet del transistor NPN).
\item Vishay / fabricante, \emph{BPW40} (datasheet del fototransistor).
\item Espressif, \emph{ESP32-S3 Series Datasheet} (GPIO y ADC).
\end{itemize}

\end{document}